% !TeX root = ../papers/nc-p3.tex
\section{異なる二つのVeronese表示を持つ階数10の正則窓}
\label{sec:twenty-point-companion}
$u,v\in k^3$に対し、標準体積形式と内積で定まる外積を$u\times v$、
$D(u,v)_i=u_jv_k+u_kv_j$（$(i,j,k)$は巡回順）とする。
$\epsilon=\pm1$について六次元への写像を
\[
 B_\epsilon((r,u),(s,v))
 =\big(rv+\epsilon u\times v,\;su+\epsilon D(u,v)\big)
\]
と置く。前節の$\beta=B_+$、$\gamma=B_-$に通常の内積を使う。

\begin{theorem}\label{thm:twenty-point-distinct-companion}
この四窓は標準14641型Koszul AS正則で反対階数10、最小周期2である。
共通点スキームは被約な二十点で、その固有な二乗像は相異なる二つの
Veronese三次元多様体に入る。同じ第一窓からの片側対称な相手は一意である。
\end{theorem}

\paragraph{点と階数。}
$B_+,B_-$を六行十六列行列とする。$w=B_+^{\mathsf t}B_-$は二つの偶位置の交換で
不変である。これは外積の内積の公式を代入しても、四添字の係数比較でも従う。
両隣接階数は6、反対切断の階数は10である。
偶文字$(r,x,y,z)$で評価した$B_+$の四次小行列式を調べる。
$r=1$の辞書式消去は
\[
 x^2-4y^2z^2-y^2-z^2,\quad 2xy^2+2xz^2+x,\quad xyz,\quad
 2y^3+y,\quad 2z^3+z
\]
が生成するイデアルを与える。台は原点と、一座標が零で残り二つが
独立に$\pm\sqrt{-1/2}$である十二点である。
$r=0$では$u\times v=0$から$v$は$u$の非零倍であり、
$s u_i+2u_ju_k=0$に規格化できる。座標軸の三点と、
三座標が符号を除き等しい四点を得る。各点で奇側の核は直線である
（$r=1,u\ne0$では$I+[u]_\times$の核が直線で$s$も決まる）。
従って六本の双一次式は有限完全交叉で、その長さ
$[h_1^3h_2^3](h_1+h_2)^6=20$と相異なる点数が等しく、全て被約である。
$J=\operatorname{diag}(-1,1,1,1)$とすると
$B_-=-B_+(J\otimes J)$で点集合も$J$に安定である。
第三・第四位置を$J$で同定すれば、反復する二窓は
$R_0=\operatorname{row}(B_+)$、$R_1=\operatorname{row}(B_+^{\rm op}(1\otimes J))$となる。
ここで$\mathrm{op}$は二文字の順序交換である。

\begin{lemma}[正規二次式を持つ二周期窓の有限認証]\label{lem:normal-parameter-six-test}
二周期の四文字六関係の窓で、各位置に四つの二次式$q_{a,p}$があり、
$q_{a,p}v=\sigma_a(v)q_{a,p+1}$が全ての文字で成立するとする。
$\sigma_a$は二周期の次数付き自己同型で全ての$q_b$を固定する。
次の条件で全次数の標準Koszul AS正則性を認証できる。
\begin{enumerate}
\item 各位置で四式を割った次元は次数0から5で$1,4,6,4,1,0$、
元の次元は次数6まで$\binom{n+3}{3}$である。
\item 三次・四次Koszul交叉の次元は4と1で、その最上テンソルは
定理\ref{thm:frobenius-tensor-reconstruction}の一文字・隣接階数を持ち巡回で対応する。
\item 四次で止めた右Koszul列$C_i=W_i\otimes A$で、
各$i=0,\ldots,4$の位置$i$、内部次数$i+2$におけるホモロジーが零である。
\end{enumerate}
\end{lemma}
\begin{proof}
始点ごとに商の斉次基底を語の線型結合$b$に持ち上げ、抽象的な重み2の
四変数多項式環$P$上でこれらを基底とする自由加群$L$を取る。
正規性で四式を右端へ移せ、互いに可換なので$b p(q)$への写像$L\to A$は全射である。
商が長さ5以上で零になることは一次生成性による。
$H_L(t)=(1+t)^4/(1-t^2)^4$なので次数6まではこの全射が同型である。
各文字の左乗法を、$b$での積をこの同型で一意に持ち上げて右$P$線型に延長する。
各二次関係の作用は$b$で評価すれば次数が高々6で零なので、行列は全ての次数で関係を満たす。
空語に作用させると$A\to L$を得る。次数2以下の一致と$P$線型性より、
$q$の任意の単項式は空語を対応する$P$係数倍へ送る。さらに各$b$も空語を
その基底へ送るため合成$L\to A\to L$は恒等である。これで全次数の自由性を得る。
二周期性により、二文字を付けた始点・終点の移動はここで同一視できる。

次に各$q_a$について$C_i$上に内部次数2の右加群鎖写像
$m_i(w\otimes b)=\sigma_a(w)\otimes q_{a,p+i} b$を入れる。
微分が末尾一文字を掛ける定義と正規性から鎖写像である。
これは係数も含めて$\sigma_a$を施した後の右$q_a$倍に等しい。
生成元上で低い$i$から順に鎖ホモトピーを作る。
既に$h_{i-1}$を選んだとき$(m_i-h_{i-1}d_i)(w\otimes1)$は
内部次数$i+2$のサイクルである。仮定3から一つ逆像を取り、
$h_i$を右線型に延ばす。$i=4$ではそのサイクルが零である。
従って右$q_a$倍は全ホモロジーを消す。
各$C_i$は有限自由$P$加群で、全ホモロジーは原点に支えられる。
長さ4のこの複体の正ホモロジーは零である。実際、原点で局所化し、
最上項から帰納すると既に完全な尾部の余核は位置$i$で深さ$i$以上になる
（自由加群の深さ4と短完全列の深さの補題）。$i>0$ではそこに有限長部分加群がない。
原点の外では初めからホモロジーが零である。
次数零のホモロジーは始点単純であり、線型解消を得る。
Koszul双対の完全対合は仮定2とこの解消から従い、左右AS性と大域次元4を与える。
\end{proof}

\begin{proof}[定理\ref{thm:twenty-point-distinct-companion}の全次数証明]
二周期座標で$ii$を位置$p$での同名二文字の積とする。
四つの$q_a$を$00,11,22,33$の係数で与える。$q_0$は両位置で$(0,1,1,1)$、
$q_1$は偶位置で$(-1,0,2,-2)$、奇位置で$(-1,0,-2,2)$とし、
$q_2,q_3$は後三座標を巡回させる。両位置でこれらは四平方の基底である。
$\sigma_0=1$、$a=1,2,3$では$\sigma_a$は$0,a$の文字を固定し残り二文字を負にする。
六関係が標数2の二次元ベクトル空間の文字次数で斉次なのでこれらは自己同型である。
三次の経路商で$ q_{a,p}v=\sigma_a(v)q_{a,p+1}$を係数ごとに検算する。

\path{computations/verify_twenty_point_companion.py}は有理数体で以下を計算する。
長さ$n$の商は、一文字と既知の長さ$n-1$の商の積から、
六関係と長さ$n-2$の商の積の列空間を引くことで作る。
四平方を追加した場合と元の場合に補題の次元列を得る。
各始点で$W_3,W_4$は次数ごとの二つの部分空間の交叉から直接計算して次元4と1になる。
$w$の巡回と基底変換は$W_4$の非零元を与え、各一文字切断の階数は4である
（軸での評価、または係数行列）。$d_i$を右係数次数2で計算した階数は
$i=1,2,3,4$の順に$20,45,36,10$、右係数次数1では$10,20,15,4$である。
隣り合う二階数の和は各内部次数$i+2$で中間項の次元と等しく、補題が適用できる。

階数10なので定理\ref{thm:bisymmetric-six-or-eight}より両対称表示を持たない。
周期1なら偶位置の対称な交換を一位置移して奇位置の対称な交換を得る。
二組は異なる因子に作用するので同時に対称化でき、矛盾する。
二周期表示と合わせて最小周期は2である。
二つのVeroneseが一致すれば命題\ref{prop:single-opposite-intrinsic-square}の証明の通り
有限グラフと二つの奇窓が一致し両対称になる。従ってこの二つは相異なる。
相手窓の一意性は命題\ref{prop:single-opposite-prolongation}である。
\end{proof}
